\section{Which binary Frobenius tests see the entire nonabelian spectrum?}
The one-sided argument also suggests an inverse question. Given the Dirichlet factors, which individual Frobenius events have a strictly negative coefficient at every remaining irreducible character? The two ProveIt reports on relative prime races and cyclic generators give useful signed tests. Imposing a binary indicator and this particular strict sign condition makes a full classification possible.

Let $H\ne1$ be a normal subgroup of a finite group $G$, $\pi:G\to G/H$, $q=[G:H]$, and $h=|H|$. For a real inversion-symmetric function $f$ put $A_{uv}=f(u^{-1}v)$. Write $V$ for the subspace of $\mathbb R^G$ consisting of vectors whose coordinates sum to zero on each coset of $H$. Simultaneous left translation by $H$ acts transitively on each coset, so the coset blocks have constant row and column sums and $A$ preserves $V$.

\paragraph{Integral rigidity.}
Suppose $f$ takes integer values and $A|_V$ is negative definite. Then
\[
 \sum_{u\in H}\bigl(f(u)-f(1)\bigr)\ge h-1.
\]
Equality is equivalent to $f=k\circ\pi-\mathbf1_{\{1\}}$ for an integer-valued, inversion-symmetric function $k$ on the quotient. In particular, for binary $f$, negative definiteness is equivalent to
\[
 f=\mathbf1_{\pi^{-1}(D)}-\mathbf1_{\{1\}},
 \qquad 1\in D=D^{-1}\subseteq G/H.
\]
Indeed $T=-A|_V$ preserves the full lattice $V\cap\mathbb Z^G$ and is positive definite, hence $\det T$ is a positive integer. On the coset-constant subspace the trace of $A$ is $q\sum_{u\in H}f(u)$, by taking normalized coset indicators as an orthonormal basis. Its full trace is $qh f(1)$. Consequently
\[
 1\le (\det T)^{1/[q(h-1)]}\le
 \frac{\operatorname{tr}T}{q(h-1)}
 =\frac{\sum_{u\in H}(f(u)-f(1))}{h-1}.
\]
Equality in the displayed inequality of the assertion forces all eigenvalues of $T$ to equal one. Then $(A+I)(e_u-e_v)=0$ for $u,v$ in the same coset. Reading row $1$ gives the quotient factorization. The reverse implication follows since the convolution matrix of a quotient pullback annihilates $V$. In the binary case the numerator is at most $h-1$, so equality is forced and $f(1)=0$.

For a conjugacy-invariant $f$ and $H=[G,G]\ne1$, $V$ is precisely the sum of the nonlinear isotypic subspaces. On the subspace of an irreducible character $\chi$, convolution has eigenvalue $|G|\langle f,\chi\rangle/\chi(1)$. Thus all strictly negative nonlinear coefficients of a binary test are forced to equal $-\chi(1)/|G|$. Its minimum support is $H\setminus\{1\}$, uniquely, of density $(h-1)/|G|$. This optimization concerns the stated uniform sign condition. Tests using virtual-character information, signed weights, or selected irreducibles have other possibilities.

\paragraph{An arithmetic consequence.}
Fix a nonabelian finite Galois extension $L/\mathbb Q$, let $G$ be its group, $H=[G,G]$, $K=L^H$, $a=[K:\mathbb Q]$, $n=[L:\mathbb Q]$, and $\delta=1/a-1/n$. At unramified primes write $g_p$ for a representative of Frobenius. Given Dirichlet GRH, the following inequality for every $\epsilon>0$ is equivalent to RH for $\zeta_L$.
\[
 \sum_{p\le x,\ g_p\in H\setminus\{1\}}\log p-\delta x
 \le C_{L,\epsilon} x^{1/2+\epsilon}.
\]
To prove this, put $f=\mathbf1_H-\mathbf1_{\{1\}}$ and
$\Psi(x)=\sum_{p^j\le x} f(g_p^j)\log p$, omitting ramified primes throughout. The terms with $j\ge2$ sum to $O(\sqrt x\log^2(2x))$. The transform of $\Psi(x)-\delta x$ against $x^{-s-1}dx$ on $[1,\infty)$ equals
\[
 T(s)=\frac1s\left(-\frac{\zeta'_K(s)}{a\zeta_K(s)}
                   +\frac{\zeta'_L(s)}{n\zeta_L(s)}+U(s)\right)
                  -\frac{\delta}{s-1}.
\]
Here $U$ is holomorphic in $\Re s>0$. Indeed the normalized unramified coefficient of $-\zeta'_F/\zeta_F$ at $p^j$, divided by $\log p$, is one precisely when the $j$th Frobenius power becomes trivial on $F$. Terms from the finitely many excluded factors are holomorphic in this half-plane. The pole at 1 cancels. Since $K$ is abelian, its Dirichlet factorization leaves no zero with $\Re s>1/2$. Consequently $T$ has a pole at each such zero $\rho$ of $\zeta_L$, with residue $m_\rho/(n\rho)$.

Choose $\theta=1/2+\epsilon<1$. The one-sided bound gives a nonnegative function $C x^\theta-\Psi(x)+\delta x$ for suitable $C$, with Mellin transform $C/(s-\theta)-T(s)$. If its finite abscissa of convergence exceeded $\theta$, Landau's lemma would force a real singularity there. The meromorphic formula permits only a simple pole of strictly negative residue at that real point; this is impossible because the transform is nonnegative as the point is approached from the right. The transform therefore converges absolutely for $\Re s>\theta$ and is holomorphic there. Varying $\theta$ excludes every zero to the right of the critical line, and the functional equation gives the assertion. Conversely, RH for the two Dedekind functions gives $\psi_F(x)=x+O_F(\sqrt x\log^2(2x))$ for both fields by the prime ideal theorem under RH \cite{gm}. Subtracting their normalized counts proves the required inequality. The contributions at omitted primes are $O_L(\log(2x))$. The other classified binary events differ by abelian Frobenius classes, whose centered sums are within the same error under Dirichlet GRH. This argument requires no Artin holomorphy conjecture.

For the splitting field of $X^3-X-1$, the group is $S_3$ by irreducibility modulo 2 and discriminant $-23$, giving $K=\mathbb Q(\sqrt{-23})$. The minimal event counts primes $p\ne23$ at which the polynomial stays irreducible, with density $1/3$. An upper square-root-scale estimate for that count remains the missing arithmetic input. Identity and the two 3-cycles have identical images in $G^{\rm ab}$, which explains why Dirichlet GRH by itself supplies no such estimate through the abelian character expansion.
